\section{Introduction}\label{intro:section}

Let $\A_2$ be the coarse moduli variety of principally polarized abelian
surfaces. A curve in $\A_2$ is \emph{Hodge generic} if it is contained in no
proper Shimura-special subvariety. The curve case of the Zilber--Pink
conjecture predicts that such a curve meets the union of all special
subvarieties of dimension at most one in only finitely many points.
The word ``all'' includes every Hecke translate; neither an isogeny degree
nor an endomorphism order is fixed.

We first prove the elliptic-square part of this prediction, then combine
it with the two companion branch theorems to obtain the full conclusion.
For a point $s\in\A_2(\Qbar)$, write $(A_s,\lambda_s)$ for its principally
polarized abelian surface, and set
\[
 \Sigma_{E^2}(C)=
 \bigl\{s\in C(\Qbar): A_s\sim E^2
       \text{ for an elliptic curve }E/\Qbar\text{ without CM}\bigr\}.
\]
Here $\sim$ denotes isogeny over $\Qbar$, without any requirement that it
respect the principal polarization. The condition is equivalent to
$\End_{\Qbar}(A_s)\otimes_{\Z}\Q\simeq M_2(\Q)$.

\begin{theorem}[The elliptic-square case]\label{intro:squares}
Let $C\subset\A_2$ be a nonempty reduced irreducible closed algebraic
curve over $\Qbar$. If $C$ is Hodge generic, then $\Sigma_{E^2}(C)$ is
finite.
\end{theorem}

No condition is imposed on the closure of $C$ at the boundary of a
compactification. In particular, the theorem includes complete curves.
The elliptic curve, its isogeny degree and the principal polarization may
all vary.

\subsection{The companion branches and the full theorem}

The two remaining special-curve loci are treated by the following
companion theorems. We state their exact conclusions to specify the
inputs to the full finiteness theorem.

\begin{theorem}[The two companion branches]\label{intro:branches}
For every nonempty reduced irreducible closed Hodge-generic curve
$C\subset\A_2$ over $\Qbar$, the following two subsets of $C(\Qbar)$
are finite:
\begin{enumerate}[label=\textup{(\Alph*)}]
 \item the points for which $A_s$ is isogenous over $\Qbar$ to
 $E_1\times E_2$, with at least one elliptic factor having CM;
 \item the points for which the full algebra
 $\End_{\Qbar}(A_s)\otimes_{\Z}\Q$ is an indefinite quaternion division
 algebra over $\Q$.
\end{enumerate}
These statements allow all CM fields, endomorphism orders, quaternion
algebras, isogeny degrees and principal polarizations. They impose no
boundary, reduction, integrality, height or field-degree condition.
\end{theorem}

In (B), indefiniteness means that the algebra becomes $M_2(\R)$ over
$\R$; the division condition excludes $M_2(\Q)$. An embedding of a
quaternion algebra into a larger endomorphism algebra does not satisfy
(B). In (A), both-CM products are included. The complete proofs of
these conclusions are given respectively in the companion articles
\cite[Theorem~1.1]{CompanionECM} and
\cite[Theorem~1.1]{CompanionQuaternion}. The proof below keeps their
use separate from the elliptic-square theorem.

\begin{theorem}[Zilber--Pink for curves in $\A_2$]
\label{intro:zp}
For every nonempty reduced irreducible
closed Hodge-generic curve $C\subset\A_2$ over $\Qbar$, the set
\[
 \Sigma(C)=C(\Qbar)\cap
 \bigcup_{\substack{Z\subset\A_2\ \mathrm{special}\mathstrut\\
                     \dim Z\le 1}} Z(\Qbar)
\]
is finite.
\end{theorem}

\Cref{intro:zp} resolves positively the Hodge-generic-curve
case of the Zilber--Pink conjecture in $\A_2$ over $\Qbar$. Its proof combines
\Cref{intro:squares,intro:branches} with the classification of special
curves and Andr\'e--Oort for special points, as explained in
\Cref{fin:classification,fin:complete}. Theorem~\ref{intro:squares}
and its proof are independent of the two companion results.
After removing special points, the three branch statements cover the
points of $C$ that are Hodge generic on a special curve;
\Cref{fin:classification}
proves this classification with the full endomorphism algebras specified.
For a curve that is not Hodge generic in $\A_2$, Andr\'e--Oort also gives
the corresponding special-closure formulation: if $S_C$ is the smallest
special subvariety containing $C$, then $C$ meets the union of special
$Z$ with $\dim Z<\dim S_C-1$ in finitely many points
(\Cref{fin:special-closure}).

\subsection{History and the arithmetic step}

Zilber's atypical-intersection principle for semiabelian varieties
\cite{Zilber2002} and Pink's formulation for mixed Shimura varieties
place this problem in a common unlikely-intersection framework.
Pink's formulation predicts finiteness for a
Hodge-generic curve meeting special subvarieties of codimension at least
two \cite[Conjecture~1.3]{Pink2005}. In the three-dimensional space
$\A_2$, the special curves have three generic types:
elliptic products with a CM factor, non-CM elliptic squares, and
quaternion-division multiplication. Pila and Tsimerman established
Andr\'e--Oort for $\A_2$, which supplies the special-point case
\cite[Theorem~1.1]{PilaTsimerman2013}. Daw and Orr developed the
corresponding unlikely-intersection arguments for these curve loci
\cite{DawOrrECM2021,DawOrr2022}.

The arithmetic-to-geometric strategy compares lower bounds coming
from Galois conjugates with upper bounds for rational points in
definable sets. Pila and Zannier introduced this strategy in their
proof of Manin--Mumford \cite{PilaZannier2008}, using the counting
theorem of Pila and Wilkie \cite{PilaWilkie2006}. In Siegel moduli,
the arithmetic difficulty is to obtain a sufficiently large orbit
uniformly as the special subvariety varies. Daw and Orr's quantitative
reduction theory supplies the geometric and counting implication
for the elliptic-square and quaternionic loci. Their unconditional
quaternionic theorem assumes that the curve meets the zero-dimensional
Baily--Borel boundary \cite[Theorem~1.4]{DawOrr2022}; the orbit-to-finiteness
implication we use has no such hypothesis.

Daw and Orr subsequently proved the full special-curve finiteness
conclusion, including elliptic squares, under the same total
multiplicative-degeneration hypothesis
\cite[Theorem~1.1 and Corollary~1.2]{DawOrrMultiplicative2025}.
Papas studied splitting points using periods centered at an interior
$E\times\mathrm{CM}$ or elliptic-square point. For an everywhere
potentially-good center, his height bound depends polynomially on the
point degree and the number of supersingular places of proximity;
his finiteness corollary fixes an upper bound for that number
\cite[Theorem~1.5 and Corollary~1.6]{PapasGFunctionsI2025}.
His analogous bad-reduction statements are conditional on comparison
conjectures \cite[Theorem~1.8 and Corollary~1.9]{PapasGFunctionsI2025}.
In the product $Y(1)^3$ of three modular curves, Daw, Orr and Papas
obtain finiteness from an interior
basepoint, allowing bad reduction but restricting the growth of the
number of supersingular places of proximity
\cite[Theorems~1.6 and~1.8]{DawOrrPapas2026}.
The present height estimate permits bad reduction of the center and
does not bound that number of places.

For the elliptic-square locus, the implication we use is Daw--Orr's
\cite[Theorem~1.3]{DawOrr2022}: a suitable polynomial lower bound for
Galois orbits in terms of the endomorphism-order discriminant implies
finiteness. Their arithmetic hypothesis is stated in
\cite[Conjecture~6.2]{DawOrr2022}. We prove the required bound for the
entire elliptic-square locus before applying that theorem. The constants
may depend on $C$; effectivity is not required.

The main arithmetic result is a polynomial height bound for selected
lifts $t$ to one fixed fine-level curve. Write
$d=\max\{2,[K(t):K]\}$, where $K$ contains the fixed data, and let
$h(t)\geq2$ be a shifted ample logarithmic height. If $s$ is the coarse
image of $t$, let $\Delta_s$ be the discriminant of the full geometric
order $\End_{\Qbar}(A_s)$, using the trace of its regular representation.
The two arithmetic estimates are
\[
 h(t)\leq c d^{\kappa},\qquad
 \Delta_s\leq c\bigl(dh(t)\bigr)^{\kappa'}.
\]
We prove the first in \Cref{nid:height}; the second comes from the
endomorphism estimates of Masser and W\"ustholz
\cite{MasserWustholz1994}, as explained in \Cref{ht:endomorphisms}.
Since the cover degree is fixed, $d\leq c[K(s):K]$. Combining the two
estimates therefore gives the required orbit bound
$[K(s):K]\geq c'\Delta_s^{\delta}$ for some $\delta>0$.
The discriminant belongs to the full order, not just the symmetric plane
used in our period equations. This is why we bound endomorphisms spanning
the whole algebra as well as a pair spanning that plane.
\Cref{fin:orbit,fin:squares} account for every fixed omission and identify
$\Delta_s$ with the complexity in the counting theorem.

Several classical inputs make this reduction quantitative. Andr\'e's
normality theorem for connected algebraic monodromy identifies the
monodromy group of our Hodge-generic family
\cite[\S5, Theorem~1]{Andre1992}. Strong approximation then permits
independent markings at two auxiliary levels \cite{MVW1984}.
Deligne's regularity theorem and Schmid's Hodge-norm estimates control
fixed algebraic frames near the punctures
\cite{Deligne1970,Schmid1973}, while rational crystalline comparison
identifies horizontal transport at nearby good reductions
\cite{BerthelotOgus1983,Berthelot1982}. The interpolation argument uses
effective polynomial ideal membership and arithmetic intersection
bounds \cite{Hermann1926,Dube1990,Aschenbrenner2004,KrickPardoSombra2001};
its graph construction is what turns those bounds into the height
estimate needed here.

The use of arithmetic auxiliary functions and period equations belongs
to the tradition of Bombieri's and Andr\'e's work on $G$-functions
\cite{Bombieri1981,Andre1989}. For curves in products of modular curves
with multiplicative degeneration, Daw and Orr use relations valid
simultaneously at all nonarchimedean places, obtained from Tate
uniformization \cite[\S1.B]{DawOrrModular2025}.
Interior $p$-adic periods are developed
in Andr\'e's motivic account \cite{Andre1995}; see also his recent
discussion \cite[\S\S2.2 and 3.3.4]{AndreSurvey2025} and
\cite[Theorem~3.4]{PapasGFunctionsI2025}. Our interpolation theorem
is formulated for fixed differential systems at an ordinary point.
\Cref{int:section} states the theorem and explains its mechanism;
\Cref{int:proof} proves the varying-curve multiplicity estimate and
two-argument graph descent.

\subsection{The proof mechanism}

Suppose the square locus is infinite, and fix one of its points $t_0$
on a smooth fine-level cover. This interior point will be the center of
our local equations. On first homology, put
\[
 \cP=\{D:D^\dagger=D,\ \Tr(D)=0\},\qquad
 \langle D,E\rangle=\Tr(DE),
\]
where $\dagger$ is the adjoint for the polarization. This is a rank-five
quadratic system. At a non-CM elliptic square, its actual rational
endomorphisms form a positive plane. Choose labeled integral pairs
$\mathbf B=(B_1,B_2)$ and $\mathbf P=(P_1,P_2)$ spanning the center and
target planes, with the target pair chosen using the endomorphism
estimate. Let $a=x(t)$ for a local parameter $x$ vanishing at $t_0$,
and let $Y(a)$ be horizontal back-transport in a fixed de Rham frame
from the target to the center, with $Y(0)=1$. We study the four cross
pairings
\[
 \langle\mathbf B,Y(a)\mathbf P\rangle
   :=\bigl(\langle B_i,Y(a)P_j\rangle\bigr)_{1\leq i,j\leq2}.
\]

When the two fibers have the same good reduction, rational cohomological
comparison identifies these entries with traces of products of actual
endomorphisms of that reduction. Frobenius preserves the traces after
simultaneously conjugating both fibers and both pairs. Writing primes
for those conjugate data gives the equation
\[
 \langle\mathbf B,Y(a)\mathbf P\rangle
   =\langle\mathbf B',Y'(aR)\mathbf P'\rangle,
 \qquad R=a'/a,\quad |R|_v=1.
\]
The residue characteristic no longer occurs in this equation. Grouping
places by the two labeled conjugate tuples gives only polynomially many
systems in $d$, even when the number of proximity primes is unbounded.
At the archimedean and finitely many exceptional places we instead use
a single equation $\langle\mathbf B,Y(a)\mathbf P\rangle=\mathbf r$
with a rational matrix $\mathbf r$.

Two fixed prime levels make these equations usable. At the first,
$\ell_1$, the separately saturated center and target plane lattices are
marked so that their reductions span a degenerate four-space; at the
second, $\ell_2$, their sum is nondegenerate. Both conditions can be
imposed on one connected finite cover. The second incidence excludes
sufficiently close targets at a place where the center is not potentially
good: both planes would preserve the same toric two-space, forcing their
sum to be degenerate. Thus every retained finite place admits the
good-reduction comparison just described.

The first incidence serves a different purpose. Full monodromy rules out
identities among the equations unless the target and its conjugate lie
on one of a fixed finite list of polarized Hecke correspondences. On that
list, isogenies of fixed multiplier transport the primed pairs back to
the unprimed fibers. Monodromy then forces any identity to identify the
transported pairs with the original pairs by a scalar on the center pair
and its reciprocal on the target pair. The center norms force that
scalar to be $+1$ or $-1$. The positive sign would
make a rational square equal to $p/m$, where $p$ is the residue
characteristic and $m$ is a fixed isogeny multiplier prime to $p$.
The negative sign contradicts the degenerate incidence at $\ell_1$.
The Hecke list and its multipliers are fixed before choosing these
levels, so $\ell_1$ can avoid all those multipliers.

There are two different size estimates in the interpolation step. Let
$U$ collect the base coordinates, the target endomorphism coordinates
and, when needed, $R$. Their global height can grow like $h=h(t)$:
$h_{\rm aff}(a,U)\leq Cd^C h$. Unless $h$ is already polynomially
bounded in $d$, we select a group of nearby places at which the weighted
positive logarithms of $U$ are at most $Cd^C(1+\log h)^C$, while
the weighted sum of $\log|a|_v^{-1}$ is at least
$h/(Cd^C)$. Here the weights are the normalized absolute-height weights.
The center divisor supplies the large proximity sum; integral lattices
preserved by endomorphisms and bounded center frames supply the much
smaller coordinate sum. These estimates hold before any bound for $h$
in terms of $d$ is known.

\Cref{int:height} turns this imbalance, together with the nonidentity
just explained, into $h\leq Cd^C$. Its proof constructs a polynomial
whose degree is small compared with its order of local smallness.
The product formula either bounds $h$ or forces the target into a
proper algebraic subvariety. The process repeats in that subvariety,
which explains why nonidentity must hold on boundary branches of
arbitrary algebraic varieties through the target, even when their
operator coordinates are not Hodge classes. \Cref{int:section} explains how a graph construction
makes this dimension descent quantitative.

The quaternionic companion develops Frobenius cross-pairing
interpolation and the degenerate auxiliary incidence
\cite[Sections~2--5]{CompanionQuaternion}. The arguments needed here, including
the second incidence, the integral coordinate bounds and the
nonidentity analysis for split endomorphism algebras, are proved below.
The interpolation theorem and the separation of bad-center disks from
good-reduction comparisons are the parts of the construction that may
be useful in other height problems.

Figure~\ref{fig:two-incidences} distinguishes the two uses of the
auxiliary levels and their roles alongside the arithmetic estimates.
\begin{figure}[t]
\centering
\begin{tikzpicture}[
  box/.style={draw=linkblue,rounded corners=2pt,align=center,
    text width=55mm,minimum height=13mm,inner sep=5pt,font=\small},
  arr/.style={-{Stealth[length=2mm]},draw=linkblue,thick},
  node distance=10mm and 8mm]
\node[box] (minus) {First level $\ell_1$\\
  center and target planes span\\a degenerate four-space};
\node[box,right=of minus] (plus) {Second level $\ell_2$\\
  center and target planes span\\a nondegenerate four-space};
\node[box,below=of minus] (sign) {Rules out the negative sign\\
  in a Frobenius identity};
\node[box,below=of plus] (torus) {Rules out a common\\
  toric two-space at a bad center};
\draw[arr] (minus)--(sign);
\draw[arr] (plus)--(torus);
\node[box,text width=70mm,below=14mm of $(sign.south)!0.5!(torus.south)$]
  (height) {Common equations and nonidentity\\
  + proximity and coordinate bounds\\
  $\Longrightarrow$ polynomial height bound};
\draw[arr] (sign.south)--(height.north west);
\draw[arr] (torus.south)--(height.north east);
\end{tikzpicture}
\caption{The two level incidences have different roles. The first is used
to exclude the negative sign in the nonidentity argument; the positive
sign is excluded by the Frobenius multiplier. The second excludes
proximity to centers with non-potentially-good reduction. The height
conclusion uses common equations at the retained places, nonidentity on
algebraic boundary branches, the large proximity sum, and the small
positive coordinate logarithms described in the text. The diagram is
schematic: all plane positions are in the separately marked finite-field
quadratic spaces.}
\label{fig:two-incidences}
\end{figure}


\subsection{Conventions and organization}

All heights are absolute logarithmic heights. Place sums use normalized
absolute-height weights, so extending a field preserves the sum for
pulled-back quantities. Constants may depend on the curve, the center,
fixed levels, frames and other fixed data. They never depend on the
varying target, its height or degree, or a varying residue characteristic.
In an estimate described as polynomial, the exponent is fixed.

We use covariant first homology and its de Rham and prime-to-residue-
characteristic \'etale realizations. The term \emph{ordinary point} for a
differential equation means a point where the connection has no
singularity; it places no ordinarity condition on the reduction of an
abelian variety. Every reduction of an endomorphism plane uses its
saturation in the specified integral lattice.

\Cref{int:section} states the interpolation theorem and explains its
proof mechanism.
\Cref{geo:section} constructs the geometric systems and controls the
exceptional Hecke correspondences. \Cref{inc:section} gives the two
incidences and excludes bad-center disks. \Cref{ht:section} proves the
coordinate estimates, and \Cref{loc:section} constructs the local
equations with their quantitative bounds. \Cref{nid:section} rules out
identities and obtains the height bound. Finally,
\Cref{fin:section} proves \Cref{intro:squares,intro:zp} and the
special-closure corollary. The complete proof of the interpolation
theorem is given in \Cref{int:proof}.
